\begin{resumo}
A adoção crescente da linguagem Go em arquiteturas modernas evidencia um desafio persistente na segurança de software: a alta taxa de falsos positivos gerados por ferramentas de Análise Estática (SAST). Esse cenário decorre da natureza inerentemente conservadora dos analisadores estáticos, que, na tentativa de não omitir ameaças reais, emitem alertas excessivos ao examinar o código, especialmente ao processar as abstrações arquiteturais e os modelos de concorrência inerentes à linguagem Go. Para mitigar a fadiga de alertas nas equipes de desenvolvimento, este trabalho tem como objetivo desenvolver uma abordagem automatizada para a triagem e filtragem de falsos positivos em relatórios de SAST, por meio do delineamento de uma arquitetura neuro-simbólica. Essa solução busca integrar a precisão estrutural do motor de análise de código com a capacidade de interpretação semântica de Modelos de Linguagem de Grande Escala (LLMs), empregados como classificadores secundários. A metodologia fundamenta-se em um experimento empírico quantitativo que acopla o motor Semgrep a modelos neurais de fronteira. O processo utiliza a extração de contexto adjacente e a sua respectiva injeção em \textit{prompts} especialistas, estruturados com regras teóricas de segurança e peculiaridades da linguagem. Como resultado, espera-se que o raciocínio guiado da IA permita filtrar o ruído analítico, reduzindo o esforço cognitivo exigido na validação manual sem comprometer a detecção de vulnerabilidades genuínas.

% O resumo deve ressaltar o objetivo, o método, os resultados e as conclusões  do documento. A ordem e a extensão  destes itens dependem do tipo de resumo (informativo ou indicativo) e do  tratamento que cada item recebe no documento original. O resumo deve ser  precedido da referência do documento, com exceção do resumo inserido no  próprio documento. (\ldots) As palavras-chave devem figurar logo abaixo do resumo, antecedidas da expressão Palavras-chave:, separadas entre si por ponto e finalizadas também por ponto. O texto pode conter no mínimo 150 e no máximo 500 palavras, é aconselhável que sejam utilizadas 200 palavras. E não se separa o texto do resumo em parágrafos.

 \vspace{\onelineskip}
    
 \noindent
 \textbf{Palavras-chave}: análise estática; linguagem Go; inteligência artificial; triagem automatizada; arquitetura neuro-simbólica; segurança de software.
\end{resumo}
